\chapter{Einleitung}

Um das Verhalten von Röntgenstrahlung beim Durchgang durch Materie zu beschreiben, werden in diesem Versuch Methoden der Dosimetrie genutzt. Hiermit wird die Wirkung von Röntgenstrahlung bei Interaktion mit Materie systematisch beschrieben und erfasst sowie ihr Vorhandensein gezeigt \cite[1]{LDdosimetrie}. Es können so die Strahlenbelastung durch eine bestimmte Röntgenquelle, d.h. die erwartete Wirkung auf Gewebe, über die Ionen-, Energie- und Äquivalenzdosen bestimmt werden, welche für den Strahlenschutz wichtige Größen sind \cite[20]{dosimetriekrieger}.

In diesem Versuch werden mithilfe eines Vollschutzröntgengeräts mit einer Kupfer (Cu)- und einer Molybdän (Mo)-Röntgenröhre als Strahlenquellen die mittlere Ionendosisleistungen der Röhren bestimmt. Für die Messungen der Röntgenstrahlung werden im Verlauf des Versuchs ein luftgefüllter Kondensator, ein Stabdosimeter und ein Geiger-Müller-Zählrohr verwendet \cite[3-5]{529va}, die jeweils Beispiele für gasgefüllte Teilchendetektoren sind \cite[179]{kolanoskiwermerdetektoren}.
Es werden zunächst die Betriebsparameter der Röntgenröhren und der Detektoren untersucht. Beim Geiger-Müller-Zählrohr wird insbesondere die charakteristische Totzeit des verwendeten Geräts bestimmt. Außerdem wird durch Messung der Zählraten des Zählrohrs hinter verschiedenen Abschirmungsmaterialien und -dicken das \uppercase{Lambert}'sches Schwächungsgesetz untersucht. Für die untersuchten Materialien wird der lineare Schwächungskoeffizient $\mu$ bestimmt. Außerdem wird durch die Messung der Dosisleistung, die vom Stabdosimeter gemessen wird, die invers quadratische Abhängigkeit der Dosis vom Abstand von der Röntgenquelle überprüft.
Zusätzlich wird als letztes noch die wichtige medizinische Anwendung von Röntgenstrahlung in Computertomogrammen betrachtet. Hier wird aus vielen 2D-Aufnahmen der von verschiedenen Bestandteilen einer (biologischen) Probe verschieden transmittierte Röntgenstrahlung eine 3D-Rekonstruktion der inneren Strukturen der Probe erstellt \cite[2]{529va}.
